\section{Orthogonality also forces divisibility by four}\label{ch14:sec:four}
The square condition of Section~\ref{ch14:sec:square} relied on the circulant pattern. The next condition holds for every Hadamard matrix, circulant or not, as long as its order is greater than $2$. Small orders show why that restriction is needed. The $1$-by-$1$ matrix $(1)$ is a Hadamard matrix of order $1$, and the matrix
\[
\begin{pmatrix}1&1\\1&-1\end{pmatrix}
\]
of Section~\ref{ch10:sec:product} is a Hadamard matrix of order $2$, but neither $1$ nor $2$ is divisible by four.

\par\noindent\begin{minipage}{\linewidth}
\begin{lemma}\label{thm:hadamardorder}
If a Hadamard matrix has order $n>2$, then $4$ divides $n$.
\end{lemma}
\end{minipage}\par

\begin{proof}
The idea is to replace the matrix by a simpler Hadamard matrix of the same order, and then to look closely at three of its rows. Two kinds of change keep a matrix Hadamard.
\begin{itemize}
\item \emph{Multiplying a column by $-1$.} This changes the sign of one entry in every row. In the dot product of two rows, the term that comes from that column is a product of two entries that have both changed sign, so it is unchanged, because $(-x)(-y)=xy$. The other terms are not touched at all. So every dot product of two rows stays the same, and the entries are still $\pm1$.
\item \emph{Rearranging the columns.} This applies the same rearrangement to every row. The dot product of two rows then consists of the same terms as before, added in a different order, so it does not change.
\end{itemize}
Neither change alters the order $n$, so it is enough to prove that $4$ divides $n$ for the new matrix. Even if the original matrix was circulant, the new one usually is not. That does not matter, because the lemma is about all Hadamard matrices, and the proof uses only two facts: the entries are $\pm1$, and different rows are orthogonal.

\emph{Step 1: make the first row consist of ones.} For each column whose entry in the first row is $-1$, multiply that column by $-1$. Afterward the first row is $(1,1,\ldots,1)$.

\emph{Step 2: the second row has as many plus signs as minus signs.} The second row is orthogonal to the first. Since the first row consists of ones, their dot product is simply the sum of the entries of the second row, so this sum is zero. If the second row has $p$ entries equal to $1$, then it has $n-p$ entries equal to $-1$, and its sum is $p-(n-p)=2p-n$. So $2p=n$. Hence $n$ is even, and the second row has $n/2$ plus signs and $n/2$ minus signs. Put $h=n/2$.

\emph{Step 3: sort the second row.} Rearrange the columns so that the $h$ columns in which the second row has a $1$ come first, and the $h$ columns in which it has a $-1$ come last. The first row still consists of ones, and the second row is now
\[
(\underbrace{1,\ldots,1}_{h},\underbrace{-1,\ldots,-1}_{h}).
\]

\emph{Step 4: look at a third row.} Since $n>2$, the matrix has a third row. Let $x$ be the sum of its first $h$ entries and $y$ the sum of its last $h$ entries. Its dot product with the first row is $x+y$. Its dot product with the second row is $x-y$, because the second row multiplies the first $h$ entries by $1$ and the last $h$ entries by $-1$. The third row is orthogonal to both, so
\[
x+y=0\qquad\text{and}\qquad x-y=0.
\]
Adding the two equations gives $2x=0$, so $x=0$. (Subtracting them gives $y=0$ as well, although we will not need it.)

\emph{Step 5: count again.} So the first $h$ entries of the third row are signs whose sum is zero. The count of Step~2 applies to them: if $q$ of these $h$ entries equal $1$, their sum is $q-(h-q)=2q-h$, and this is zero, so $h=2q$. Therefore $n=2h=4q$, and $4$ divides $n$.
\end{proof}

To see the proof in action, apply it to the circulant Hadamard matrix $H$ of Section~\ref{ch14:sec:circulant}. Its first row is $(1,1,1,-1)$, so Step~1 multiplies the last column by $-1$:
\[
\begin{pmatrix}
1&1&1&-1\\
-1&1&1&1\\
1&-1&1&1\\
1&1&-1&1
\end{pmatrix}
\quad\longrightarrow\quad
\begin{pmatrix}
1&1&1&1\\
-1&1&1&-1\\
1&-1&1&-1\\
1&1&-1&-1
\end{pmatrix}.
\]
The second row is now $(-1,1,1,-1)$, with its plus signs in columns $1$ and $2$ (counting from $0$). Step~3 moves these two columns to the front, so the new column order is $1,2,0,3$:
\[
\begin{pmatrix}
1&1&1&1\\
1&1&-1&-1\\
-1&1&1&-1\\
1&-1&1&-1
\end{pmatrix}.
\]
The first row consists of ones, and the second row is $(1,1,-1,-1)$. The third row is $(-1,1,1,-1)$. Its first half $(-1,1)$ and its second half $(1,-1)$ both add up to zero, as Step~4 predicts. Each half contains one plus sign, so $h=2\cdot1$ is even and $n=4\cdot1$ is divisible by four. You can check that the rows of the final matrix are still orthogonal. Notice also that it is no longer circulant: its second row is not its first row moved one place to the right.

Now we can combine the two necessary conditions.

\par\noindent\begin{minipage}{\linewidth}
\begin{proposition}\label{thm:circulantnecessary}
If a circulant Hadamard matrix has order $n$, then $n=1$ or $n=4u^2$ for some positive integer $u$.
\end{proposition}
\end{minipage}\par

\begin{proof}
Let $a$ be the first row of the matrix. The matrix is the circulant matrix of $a$, so by Proposition~\ref{ch14:prop:rows}(b), $C_s(a)=0$ for every shift $s$ with $1\le s\le n-1$. Equation~\eqref{eq:squareorder} therefore gives $n=r^2$, where $r$ is the row sum, an integer.

If $n=1$, there is nothing more to prove. Otherwise $n\ge2$. Moreover $n\ne2$, because $2$ is not the square of an integer: $1^2=1<2<4=2^2$. Hence $n>2$, and Lemma~\ref{thm:hadamardorder} shows that $4$ divides $n=r^2$. In particular $r^2$ is even, and so $r$ is even. Indeed, if $r$ were odd, say $r=2k+1$ for an integer $k$, then $r^2=4k^2+4k+1$ would be odd. Write $|r|=2u$, where $u$ is a nonnegative integer. Since $n=r^2$ is positive, $r\ne0$, and so $u\ge1$. Then $n=r^2=|r|^2=4u^2$.
\end{proof}

The proposition allows the orders $1$ and $4u^2$ for $u=1,2,3,\ldots$, that is,
\[
1,\ 4,\ 16,\ 36,\ 64,\ 100,\ \ldots.
\]
For the first two there are examples: the $1$-by-$1$ matrix $(1)$, and the matrix $H$ of order four. For the orders $16,36,64,\ldots$ the proposition says nothing either way. It neither builds a matrix of such an order nor rules one out. This is typical of a necessary condition: it removes most orders from consideration and tells us precisely which ones remain open. That list of remaining orders is the boundary between what our elementary arguments can decide and what needs new ideas.
